
\section{Physical conditions of renormalisation}

\SourcePage{538}{543}

The theory presented so far in this chapter has been to a significant extent formal in character. We operated with all quantities as if they were finite, intentionally leaving aside the infinities arising in the theory. However, in an actual computation of the functions~$\mathcal{D}$, $\mathcal{G}$, $\Gamma$ by perturbation theory, divergent integrals appear which cannot, without additional considerations, be ascribed any definite value. The emergence of such divergences is a manifestation of the logical incompleteness of quantum electrodynamics. We shall see, however, that in this theory one can establish certain prescriptions that allow an unambiguous ``subtraction of infinities'' and the obtaining of finite values for all quantities possessing direct physical meaning. At the basis of these prescriptions lie obvious physical requirements: that the photon mass be zero and that the charge and mass of the electron equal their observed values.

Let us begin with the conditions on the photon propagator.

Consider scattering through single-particle intermediate states containing one virtual photon. The amplitude must have a pole at~$P^2 = 0$ (the square of the total 4-momentum of the initial particles~$P$ coincides with the square of the real photon mass). As we saw in~\S79, this follows from the unitarity condition. The pole term arises from diagrams of the type~(79.1):

% [Feynman diagram (110.1): scattering via single virtual photon --- two shaded blobs (representing exact scattering amplitudes) connected by a bold dashed line (exact photon propagator) carrying momentum k = P]
\begin{equation}
\label{eq:110.1}
\end{equation}

\noindent with radiative corrections, both parts of the diagram being connected by a bold dashed line (exact photon propagator). Therefore~$\mathcal{D}(k^2)$ must have a pole at~$k^2 = 0$:
\begin{equation}
\mathcal{D} \to \frac{4\pi Z}{k^2} \quad \text{when} \quad k^2 \to 0,
\label{eq:110.2}
\end{equation}
where~$Z$ is a constant. For the polarisation operator this means, according to~(103.21):
\begin{equation}
\mathcal{P}(0) = 0.
\label{eq:110.3}
\end{equation}
The coefficient~$1/Z = 1 - \mathcal{P}(k^2)/k^2\big|_{k^2 \to 0}$.

Further restrictions on~$\mathcal{P}(k^2)$ follow from the physical definition of electric charge. Two classical (arbitrarily heavy) particles at rest at a large distance from one another must interact according to the Coulomb law~$U = e^2/r$ ($r \gg 1/m$). On the other hand, this interaction is expressed by the diagram:

% [Feynman diagram (110.4): scattering of two heavy classical particles via a single virtual photon --- upper and lower straight lines (heavy particles a,b and c,d) connected by a bold dashed photon line carrying momentum k]
\begin{equation}
\label{eq:110.4}
\end{equation}

\noindent where the upper and lower lines correspond to the classical particles. The photon self-energy corrections are included in the virtual photon line. Other corrections, involving the heavy-particle lines, vanish in the limit~$M \to \infty$: adding internal lines to the diagram~(\ref{eq:110.4}) (for example, connecting lines~$a$ and~$c$, or~$a$ and~$b$, by a photon line) leads to virtual heavy-particle propagators containing the mass~$M$ in the denominator and vanishing as~$M \to \infty$.

\SourcePage{539}{544}

From the diagram~(\ref{eq:110.4}), the factor~$e^2\mathcal{D}(k^2)$ must represent (up to sign) the Fourier transform of the interaction potential. Stationarity of the particles means the virtual photon frequency~$\omega = 0$; large distances mean small wave vectors~$\mathbf{k}$. The Fourier transform of the Coulomb potential is~$4\pi e^2/\mathbf{k}^2$. Since~$\mathcal{D}$ depends only on~$k^2 = \omega^2 - \mathbf{k}^2$, we obtain:
\begin{equation}
\mathcal{D} \to \frac{4\pi}{k^2}, \quad k^2 \to 0.
\label{eq:110.5}
\end{equation}
Thus the coefficient~$Z$ in~(\ref{eq:110.2}) must equal unity: $Z = 1$ (the sign in~(\ref{eq:110.5}) is obvious: $\mathcal{D}(k^2)$ tends to the free photon propagator~$D(k^2)$). For the polarisation operator this means:
\begin{equation}
\frac{\mathcal{P}(k^2)}{k^2} \to 0, \quad k^2 \to 0.
\label{eq:110.6}
\end{equation}
Besides condition~(\ref{eq:110.3}), this also gives:
\begin{equation}
\mathcal{P}'(0) = 0.
\label{eq:110.7}
\end{equation}

In~\S103 it was noted that the effective external photon line factor~(103.15), with account of~(103.16) and~(103.20), is:
\[
\Bigl[1 + \frac{\mathcal{P}(0)\,\mathcal{D}(0)}{4\pi}\Bigr]\,e^\mu.
\]
Now, with~(\ref{eq:110.5}) and~(\ref{eq:110.6}), the correction term vanishes. In other words, in external photon lines one need not include radiative corrections at all.

Thus natural physical requirements lead to the establishment of definite (zero) values of~$\mathcal{P}(0)$ and~$\mathcal{P}'(0)$. Computing these quantities from perturbation-theory diagrams would lead to divergent integrals. The method of eliminating infinities consists in ascribing to the divergent expressions the values prescribed by the physical requirements. This procedure is called \emph{renormalisation} of the corresponding quantities.\footnote{The idea of this approach was first expressed by Kramers (H.~Kramers, 1947). Systematic application of the renormalisation method in quantum electrodynamics was carried out in the works of Dyson, Tomonaga (S.~Tomonaga), Feynman and Schwinger.}

\SourcePage{540}{545}

The procedure can also be formulated differently. For charge renormalisation one introduces a non-physical ``bare'' charge~$e_0$ as the parameter in the electromagnetic interaction Hamiltonian in the formal perturbation theory. After this, the renormalisation condition is formulated as: $e_0^2\,\mathcal{D}(k^2) \to 4\pi e^2/k^2$ (at~$k^2 \to 0$), where~$e$ is the true physical charge. From this~$e_0^2\,Z = e^2$, and together with~$Z$ the non-physical quantity~$e_0$ is eliminated. By requiring~$Z = 1$ from the outset, we perform the renormalisation ``on the fly'' and avoid introducing fictitious quantities even in intermediate calculations.

Now let us turn to the conditions on the electron propagator. Consider scattering through a single-particle intermediate state containing one virtual electron. The amplitude must have a pole when the square of the total 4-momentum of the initial particles~$P_i$ coincides with the square of the real electron mass: $P_i^2 = m^2$. The pole term in the amplitude arises from diagrams of the type

% [Feynman diagram (110.8): scattering via single virtual electron --- two blobs connected by a bold solid line (exact electron propagator) carrying momentum p = P_i = P_f, with external lines P_f on left and P_i on right]
\begin{equation}
\label{eq:110.8}
\end{equation}

\noindent where, with radiative corrections taken into account, the bold solid line is the exact electron propagator. This means that the function~$\mathcal{G}(p)$ must have a pole at~$p^2 = m^2$, i.e.\ must have the limiting form
\begin{equation}
\mathcal{G}(p) \approx Z_1\,\frac{\gamma p + m}{p^2 - m^2 + i0} + g(p), \qquad p^2 \to m^2,
\label{eq:110.9}
\end{equation}
where~$Z_1$ is a scalar constant, and~$g(p)$ remains finite as~$p^2 \to m^2$. The matrix structure of the pole term in~(\ref{eq:110.9}) (proportionality to~$\gamma p + m$) is a consequence of the same unitarity condition from which the requirement of the existence of the pole itself arose. Let us demonstrate this, simultaneously clarifying an important question about the renormalisation conditions for external electron lines.

If~$\mathcal{G}(P)$ has the limiting form~(\ref{eq:110.9}), the inverse matrix

\SourcePage{541}{546}

\begin{equation}
\mathcal{G}^{-1}(p) \approx \frac{1}{Z_1}\,(\gamma p - m) - (\gamma p - m)\,g(\gamma p - m), \qquad p^2 \to m^2.
\label{eq:110.10}
\end{equation}
The mass operator
\begin{equation}
\mathcal{M} = G^{-1} - \mathcal{G}^{-1} \approx \Bigl(1 - \frac{1}{Z_1}\Bigr)(\gamma p - m) + (\gamma p - m)\,g(\gamma p - m), \qquad p^2 \to m^2.
\label{eq:110.11}
\end{equation}

The effective external (say, incoming) electron line carries the factor (cf.~(103.15))
\begin{equation}
\mathcal{U}(p) = u(p) + \mathcal{G}(p)\,\mathcal{M}(p)\,u(p),
\label{eq:110.12}
\end{equation}
where~$u(p)$ is the ordinary spinor amplitude of the electron, satisfying the Dirac equation $(\gamma p - m)u = 0$. By virtue of the requirements of relativistic invariance ($\mathcal{U}$, like~$u$, is a bispinor), the limiting value of~$\mathcal{U}(p)$ as~$p^2 \to m^2$ can differ from~$u(p)$ only by a constant scalar factor:
\begin{equation}
\mathcal{U}(p) = Z'\,u(p).
\label{eq:110.13}
\end{equation}

This factor~$Z'$ is connected in a definite way with~$Z_1$, but finding this connection by direct substitution of~(\ref{eq:110.10}),~(\ref{eq:110.11}) into~(\ref{eq:110.12}) is not possible, owing to the indeterminacy that arises: the result depends on the order in which the limit is taken for the various factors in~(\ref{eq:110.12}).

\SourcePage{542}{547}

One can, however, avoid the question about the correct way to take the limit, by appealing instead to the unitarity condition applied to the reaction depicted by diagram~(\ref{eq:110.8}). The unitarity relation applies, generally speaking, not to individual diagrams but to amplitudes of processes as a whole. But when~$p^2 \to m^2$ the pole diagram~(\ref{eq:110.8}) gives the dominant contribution to the corresponding amplitude~$M_{fi}$, so that the other diagrams belonging to the same reaction can be disregarded.

By virtue of the unitarity requirements, as was shown in~\S\,79, the single-particle intermediate state leads to the appearance in the imaginary part of the reaction amplitude of a term with a $\delta$-function:
\begin{equation}
i\pi\delta(p^2 - m^2)\sum_{\text{pol}} M_{fn}\,M_{in}^*,
\label{eq:110.14}
\end{equation}
where in the present case the index~$n$ refers to a state containing one real electron, and the summation over polarisations is carried out as required by the symmetrisation of both sides of the unitarity relation over the helicities of the initial and final particles (as in~\S\,79; then~$M_{fi} = M_{if}$). The amplitude~$M_{fn}$ corresponds to the process depicted by the diagram

% [Feynman diagram: scattering amplitude M_{fn} --- blob on the left with outgoing lines P_f, connected to a bold solid line carrying momentum p, ending at a vertex (black dot) on the right]

\noindent and has the form
\[
M_{fn} = (M'_{fn}\,\mathcal{U}) = Z'(M'_{fn}\,u),
\]
where~$M'_{fn}$ is a quantity with one free bispinor index.\footnote{Strictly speaking, a clarification is needed here. An electron as a stable particle cannot actually decay into a different set of real particles. One can, however, formally regard as ``final particles'' certain fictitious particles with masses that would permit such a decay. The resulting relation should then be understood in the sense of analytic continuation to the real values of the masses.} Analogously the amplitude~$M_{in}^*$ has the structure
\[
M_{in}^* = (\overline{\mathcal{U}}\,M'^*_{in}) = Z'(\bar{u}\,M'^*_{in}).
\]

\SourcePage{543}{548}

The summation over electron polarisations replaces the product $(M'_{fn}\,u)(\bar{u}\,M'^*_{in})$ by~$M'_{fn}(\gamma p + m)\,M'^*_{in}$, so that the term~(\ref{eq:110.14}) in the amplitude~$M'_{fi}$ takes the form
\[
Z'^2\,i\pi\delta(p^2 - m^2)\bigl[M'_{fn}(\gamma p + m)\,M'^{*}_{in}\bigr].
\]

On the other hand, computing this same amplitude directly from diagram~(\ref{eq:110.8}) gives
\[
iM_{fi} = iM'_{fn} \cdot i\mathcal{G}(p) \cdot iM'^{*}_{in}.
\]
Comparing the two expressions confirms the limiting form for~$\mathcal{G}(p)$ (the first term in~(\ref{eq:110.9})), with
\begin{equation}
Z' = \sqrt{Z_1}.
\label{eq:110.15}
\end{equation}

Let us now show that once the required limiting form of the electron propagator has been established, there is no need for any further conditions on the vertex operator.

Consider the diagram

% [Feynman diagram (110.16): electron scattering in an external field A^{(e)}(k) --- incoming electron line (p_1) to exact vertex (black dot), outgoing electron line (p_2), with external photon (k) at the vertex]
\begin{equation}
\label{eq:110.16}
\end{equation}

\noindent describing electron scattering in an external field~$A^{(e)}(k)$ (in the first Born approximation), with all radiative corrections taken into account. In the limit~$k \to 0$, $p_2 \to p_1 \equiv p$, the self-energy corrections to the external field line vanish (recall that these corrections vanish for any~$k$ when~$k^2 = 0$). The diagram then corresponds to the amplitude
\begin{equation}
M_{fi} = -e\overline{\mathcal{U}}(p)\,\Gamma(p,p;\,0)\,\mathcal{U}(p)\,A^{(e)}(k \to 0)
\label{eq:110.17}
\end{equation}
--- the product of~$A^{(e)}$ with the electron transition current~$\overline{\mathcal{U}}\,\Gamma\,\mathcal{U}$. But when~$k \to 0$ the potential~$A^{(e)}(x)$ reduces to one independent of coordinates and time. Such a potential does not correspond to any physical field (it is a particular case of a gauge transformation), so that it cannot cause any change in the electron current. In other words, in the limit under consideration the transition current~$\overline{\mathcal{U}}\,\Gamma\,\mathcal{U}$ must simply coincide with the free current~$\bar{u}\gamma^\mu u$:
\begin{equation}
\overline{\mathcal{U}}(p)\,\Gamma^\mu(p,p;\,0)\,\mathcal{U}(p) = Z_1\,\bar{u}(p)\,\Gamma^\mu\,u(p) = \bar{u}(p)\,\gamma^\mu\,u(p).
\label{eq:110.18}
\end{equation}

\SourcePage{544}{549}

This is, in essence, also an expression of the definition of the physical charge of the electron. It is easy to see that it is automatically satisfied regardless of the value of~$Z_1$. Indeed, substituting~$\mathcal{G}^{-1}(p)$ from~(\ref{eq:110.10}) into the Ward identity~(\ref{eq:108.8}), we find
\[
\Gamma^\mu(p,p;\,0) = Z_1^{-1}\,\gamma^\mu - \gamma^\mu\,g(p)\,(\gamma p - m) - (\gamma p - m)\,g(p)\,\gamma^\mu,
\]
and equality~(\ref{eq:110.18}) is satisfied by virtue of the equations $(\gamma p - m)u = 0$, $\bar{u}(\gamma p - m) = 0$.

We see that in constructing the amplitude of a physical process the ``renormalisation constant''~$Z_1$ drops out altogether. Moreover, taking advantage of the indeterminacy arising from the divergences in computing~$\Gamma$, one can simply require that
\begin{equation}
\bar{u}(p)\,\Gamma^\mu(p,p;\,0)\,u(p) = \bar{u}(p)\,\gamma^\mu\,u(p), \qquad p^2 = m^2,
\label{eq:110.19}
\end{equation}
i.e.\ set~$Z_1 = 1$. The convenience of this definition is that the need for introducing corrections to external electron lines disappears: we simply have
\[
\mathcal{U}(p) = u(p).
\]
One can also verify this directly: noting that when~$Z_1 = 1$ the mass operator~(\ref{eq:110.11}) becomes
\begin{equation}
\mathcal{M} = (\gamma p - m)\,g\,(\gamma p - m)
\label{eq:110.20}
\end{equation}
and the second term in~(\ref{eq:110.12}) obviously vanishes. Thus, external lines --- of both photons and electrons --- will not require ``renormalisation'' corrections.\footnote{In renormalising the photon propagator the condition~$Z = 1$ arose as a necessary physical requirement, and after this the disappearance of corrections to external photon lines occurs automatically. From a formal viewpoint, however, the situation for photon and electron external lines is analogous: if~$Z \neq 1$ then the photon wave amplitude~$e_\mu$ of a real photon would be multiplied by~$\sqrt{Z}$.}
